\chapter{ORIUS as a Universal Safety Layer}
\label{ch:universal-runtime-layer}

ORIUS is organized as a typed runtime contract rather than as a monolithic controller.
The contract takes in raw telemetry, evaluates observation reliability, inflates the
uncertainty set around the state that could be physically true, tightens the admissible
action set, repairs or replaces the candidate action, and emits a certificate that can
be audited downstream.  This is the Detect-Calibrate-Constrain-Shield-Certify kernel.

The critical design choice is adapterization.  Domain logic enters the framework
through a narrow interface: telemetry parsing, uncertainty semantics, action feasibility,
repair projection, and certificate payload construction.  Everything else belongs to the
universal safety layer itself.  That is why the monograph can treat batteries,
vehicles, industrial systems, healthcare monitoring, navigation, and aerospace as
separate domain chapters while still defending one architectural object.

\paragraph{Adapter contract}\label{par:monograph-adapter-contract}
The adapter contract is the single runtime boundary that every domain must satisfy.
It is intentionally narrow so that all domain-specific logic is isolated from the
universal kernel while still producing comparable benchmark and certificate surfaces.

The layer is intentionally post-nominal.  ORIUS does not assume ownership of the
upstream forecaster, planner, or optimizer.  It wraps them.  This makes the framework
more realistic for operational deployment because most real systems inherit legacy
nominal controllers that cannot be rewritten from scratch simply to accommodate a new
safety argument.
